% Package loads shared by the whole book. Keep hyperref last.

\usepackage{geometry}
\geometry{a4paper, inner=30mm, outer=25mm, top=25mm, bottom=25mm}

\usepackage{microtype}

% Mathematics. Shapley axioms, the Integrated Gradients path integral and the
% metric definitions of Part IV are worked in the running text, so amsthm's
% environments (set up in macros.tex) and amsmath are load-bearing.
\usepackage{amsmath}
\usepackage{amssymb}
\usepackage{amsthm}

\usepackage{booktabs}
% longtable is here for two tables that grow with the chapters and cannot be
% page-broken as plain tabulars: appendix B's keep-in-English block and
% appendix D's 32-row corpus table. The rest of the book's tables are short,
% fixed-size and stay tabular; do not convert them.
\usepackage{longtable}
\usepackage{enumitem}

% csquotes ships no style for Vietnamese. Modern Vietnamese typography uses
% the same curly double quotes as English, so aliasing the two is the correct
% mapping rather than a workaround; without it \enquote{} falls back to the
% package default and warns on every run. csquotes must stay before biblatex,
% which detects it and adjusts its own quoting.
\usepackage{csquotes}
\DeclareQuoteAlias{english}{vietnamese}

\usepackage{graphicx}
\graphicspath{{figures/images/}{figures/tikz/}}

\usepackage{tikz}
\usetikzlibrary{arrows.meta, positioning, calc, fit}
% TikZ externalization caches compiled figures between runs; enable it when
% figure-heavy chapters slow the build down:
%   1. add -shell-escape to $lualatex in .latexmkrc
%   2. uncomment the two lines below
% \usetikzlibrary{external}
% \tikzexternalize[prefix=tikz-cache/]

% Pseudo-code. This book ships no code that runs anywhere: every algorithm
% (LIME's sampling loop, KernelSHAP, a perturbation metric) is typeset here
% as pseudo-code, so a reader is never asked to believe an invented listing
% is real code. There is deliberately no minted load; admitting one is a
% decision-log row in SPEC.md, not an edit here.
\usepackage{algorithm}
\usepackage{algpseudocode}

% The end-of-chapter summary box (tomtat, defined in macros.tex). breakable
% matters: a summary can straddle a page break.
\usepackage[most]{tcolorbox}
\tcbuselibrary{breakable}

\usepackage[style=numeric, backend=biber]{biblatex}
% biblatex ships no vietnamese.lbx, so it warns and falls back once the main
% language is Vietnamese. Map it to English on purpose: every entry in
% refs.bib is an English-language paper, so an English "In:" and "visited on"
% is what the entries should read anyway.
\DeclareLanguageMapping{vietnamese}{english}

% biblatex will not break a URL anywhere except after a slash, so one long
% documentation or engineering-blog URL runs past the measure and the build
% gate fails on an overfull box in a bibliography nobody edited. These permit a
% break after a digit, an uppercase letter and a lowercase letter respectively.
\setcounter{biburlnumpenalty}{9000}
\setcounter{biburlucpenalty}{9000}
\setcounter{biburllcpenalty}{9000}

\usepackage{imakeidx}

\usepackage[hidelinks]{hyperref}
